\section{Canonical high cells and endpoint transport}
\label{sec:canonical-high-cells-v3}

This section estimates the contribution of the large overlap cells. The finite
combinatorics and the phase asymptotics are kept separate. We first sum the
entire physical matching fiber for a fixed high-cell demand. We then compare
realized multiplicities with full containment, sum all positive deficits, and
regroup the full-containment references by endpoint table. Only in the final
subsection do we insert the four-size phase estimates. No residual local reward
or cycle-space factor is charged here; those factors belong to Section~9.

\subsection{The complete physical fiber}

Let $U$ be the largest class size in the fixed four-size profile and put
\[
  R_0:=\left\lfloor\frac U2\right\rfloor.
\]
A cell is \emph{high} if its multiplicity exceeds $R_0$. Two high cells cannot
share a row or a column: otherwise that row or column would contain more than
$U$ stubs. Hence all high cells form a canonical bipartite matching, denoted by
$P$.

For $e\in P$, let $s_e$ and $t_e$ be its endpoint block sizes and set
\[
  m_e:=\min\{s_e,t_e\},
  \qquad
  d_e:=|s_e-t_e|,
  \qquad
  j_e:=m_e-h_e.
\]
Thus $m_e$ is the full-containment multiplicity, $j_e$ is the realized
multiplicity, and $h_e$ is the deficit. Since $j_e>U/2$ and $m_e\le U$,
\begin{equation}
  2h_e<m_e.
  \label{eq:half-deficit-v3}
\end{equation}
Put $J:=\sum_{e\in P}j_e$.

A size-$j$ partial matching between labeled blocks of sizes $s$ and $t$ can be
chosen in
\[
  \binom sj\binom tj j!
  =
  \frac{(s)_j(t)_j}{j!}
\]
ways: choose its two endpoint sets and then the bijection between them.

\begin{proposition}[Completion-free aggregate weight]
\label{prop:aggregate-high-weight-v3}
For a fixed matching support $P$ and multiplicity vector
$j=(j_e)_{e\in P}$, the sum of the bare weights of all physical partial
matchings realizing these data is
\begin{equation}
  w(P,j)
  =
  \frac{
    \displaystyle\prod_{e\in P}(s_e)_{j_e}(t_e)_{j_e}
  }{
    \displaystyle(n)_J\prod_{e\in P}j_e!
  }
  \prod_{e\in P}g(j_e).
  \label{eq:aggregate-high-weight-v3}
\end{equation}
\end{proposition}

\begin{proof}
Because $P$ is a matching, distinct selected cells use disjoint row and column
stub sets. The complete physical fiber is therefore the Cartesian product of
the one-cell fibers, and its cardinality is the product of the counts above.
A prescribed set of $J$ stub pairs occurs with probability $(n)_J^{-1}$ in the
uniform bipartite configuration matching. Multiplying this probability by the
fiber cardinality and by the exposed reward
$\prod_{e\in P}g(j_e)$ gives \eqref{eq:aggregate-high-weight-v3}.
\end{proof}

\begin{remark}
The matching hypothesis is load-bearing. If two selected cells shared a row or
a column, their endpoint choices would compete for the same stubs and the
one-cell fiber counts would not multiply.
\end{remark}

\subsection{Deficits and the single ambient loss}

For endpoint sizes $m$ and $m+d$, define
\begin{equation}
  R_{m,d}(h)
  :=
  \frac{\binom mh}
       {(d+1)(d+2)\cdots(d+h)}
  2^{-hm+h(h+1)/2},
  \qquad
  R_{m,d}(0):=1.
  \label{eq:local-deficit-ratio-v3}
\end{equation}
The product in the denominator is empty when $h=0$.

\begin{lemma}[Exact one-cell deficit ratio]
\label{lem:exact-local-ratio-v3}
The ratio of the aggregate one-cell factor at multiplicity $m-h$ to its
full-containment value at multiplicity $m$ is $R_{m,d}(h)$.
\end{lemma}

\begin{proof}
The ratio of the physical matching counts is
\[
  \frac{(m)_{m-h}(m+d)_{m-h}}{(m-h)!}
  \frac{m!}{(m)_m(m+d)_m}
  =
  \frac{\binom mh}{(d+1)(d+2)\cdots(d+h)}.
\]
The local signed rewards satisfy
\[
  \frac{g(m-h)}{g(m)}
  =
  2^{-hm+h(h+1)/2}.
\]
Multiplying the two identities proves the claim.
\end{proof}

Write
\[
  w_{\mathrm{full}}(P):=w(P,(m_e)_{e\in P}),
  \qquad
  J_*:=\sum_{e\in P}m_e,
  \qquad
  H:=J_*-J=\sum_{e\in P}h_e.
\]
After the one-cell ratios have been extracted, the only factor that still
couples different cells is the ambient denominator:
\begin{equation}
  \frac{(n)_{J_*}}{(n)_J}
  =(n-J)_H
  \le n^H.
  \label{eq:one-global-falling-factorial-v3}
\end{equation}

\begin{lemma}[Aggregate deficit comparison]
\label{lem:aggregate-deficit-comparison-v3}
For every admissible deficit vector $h=(h_e)_{e\in P}$,
\begin{equation}
  w(P,m-h)
  \le
  w_{\mathrm{full}}(P)
  \prod_{e\in P}n^{h_e}R_{m_e,d_e}(h_e).
  \label{eq:aggregate-deficit-comparison-v3}
\end{equation}
\end{lemma}

\begin{proof}
Apply Lemma~\ref{lem:exact-local-ratio-v3} in each selected cell, apply
\eqref{eq:one-global-falling-factorial-v3} once, and use
$n^H=\prod_{e\in P}n^{h_e}$.
\end{proof}

The order of operations matters: the physical fiber is summed first, and the
finite-population loss is charged only once afterward.

\subsection{Summing all positive deficits}

For a selected cell $e$, enlarge the exact positive-deficit range to
\[
  A_e:=\{h\in\mathbb N:h\ge1,\ 2h<m_e\}.
\]
All weights are nonnegative, so the enlargement preserves the direction of the
upper bound.

\begin{lemma}[Optional-choice product]
\label{lem:optional-choice-product-v3}
Let $u_e(h)\ge0$ for $e\in P$ and $h\in A_e$. Then
\[
  \sum_{\omega}
    \prod_{e\in P}u_e(\omega_e)
  =
  \prod_{e\in P}
    \left(1+\sum_{h\in A_e}u_e(h)\right),
\]
where each $\omega_e$ is either the zero-deficit choice, of weight one, or one
positive deficit $h\in A_e$.
\end{lemma}

\begin{proof}
Expand the finite product. Each monomial records exactly one independent choice
in every selected cell, and every optional-choice vector occurs once.
\end{proof}

For $2h<m$, integrality gives $2h+1\le m$. Hence
\[
  \frac{h+1}{2}
  \le
  \frac{m+1}{4}
  \le
  m-\left\lfloor\frac{3m-1}{4}\right\rfloor.
\]
Multiplying by $h$ yields
\begin{equation}
  h\left\lfloor\frac{3m-1}{4}\right\rfloor
  \le
  hm-\frac{h(h+1)}2.
  \label{eq:three-quarter-budget-v3}
\end{equation}
Together with $\binom mh\le m^h$ and
\eqref{eq:local-deficit-ratio-v3}, this gives
\begin{equation}
  n^hR_{m,d}(h)
  \le
  \left(
    \frac{nm}{2^{\lfloor(3m-1)/4\rfloor}}
  \right)^h.
  \label{eq:three-quarter-charge-v3}
\end{equation}

Index the four endpoint sizes by $i,j\in\{0,1,2,3\}$, let $m_{ij}$ be the
smaller endpoint size of type $(i,j)$, and put
\[
  \rho_{ij}
  :=
  \frac{n m_{ij}}{2^{\lfloor(3m_{ij}-1)/4\rfloor}},
  \qquad
  \rho_{16}:=\sum_{i,j=0}^{3}\rho_{ij}.
\]
Every local base in \eqref{eq:three-quarter-charge-v3} is at most
$\rho_{16}$.

\begin{lemma}[Coarse local deficit sum]
\label{lem:coarse-local-deficit-sum-v3}
Suppose $\rho_{16}\le1$. Then, for every selected cell $e$,
\[
  \sum_{h\in A_e}n^hR_{m_e,d_e}(h)
  \le
  (\alpha+1)\rho_{16}.
\]
\end{lemma}

\begin{proof}
If $e$ has endpoint type $(i,j)$, then
\eqref{eq:three-quarter-charge-v3} bounds its term of deficit $h$ by
$\rho_{ij}^h$. Since $h\ge1$ and
$\rho_{ij}\le\rho_{16}\le1$, this is at most $\rho_{16}$. Moreover,
$|A_e|\le m_e\le\alpha+1$. Summation proves the bound.
\end{proof}

\begin{proposition}[Fixed-support all-deficit bound]
\label{prop:fixed-support-all-deficit-v3}
If $\rho_{16}\le1$, then
\begin{equation}
  \sum_{h:\,m-h\text{ high}}w(P,m-h)
  \le
  w_{\mathrm{full}}(P)
  \left(1+(\alpha+1)\rho_{16}\right)^{|P|}.
  \label{eq:fixed-support-all-deficit-v3}
\end{equation}
\end{proposition}

\begin{proof}
Use Lemma~\ref{lem:aggregate-deficit-comparison-v3}, enlarge the deficit range
to $A_e$, and apply Lemmas~\ref{lem:optional-choice-product-v3} and
\ref{lem:coarse-local-deficit-sum-v3}.
\end{proof}

Projection onto the row endpoint is injective on a block matching. Therefore,
if $K=\sum_{i=0}^{3}k_i$ is the number of row-class slots, then
\begin{equation}
  |P|\le K.
  \label{eq:support-cardinality-v3}
\end{equation}

\subsection{Regrouping by endpoint table}

For a realized endpoint table $L$, define
\[
  W(L)
  :=
  \sum_{P:\,L(P)=L}w_{\mathrm{full}}(P).
\]
Equivalently, $W(L)$ is the sum of the full-containment atoms over all
physically decorated supports with table $L$. A realized table belongs to the
finite image of the support map; it is not an arbitrary element of
$\mathbb N^{4\times4}$.

\begin{proposition}[Reference grouping]
\label{prop:reference-grouping-v3}
The full-containment references satisfy
\begin{equation}
  \sum_P w_{\mathrm{full}}(P)
  =
  \sum_{L\ \mathrm{realized}}W(L).
  \label{eq:reference-grouping-v3}
\end{equation}
More generally, for every nonnegative function $F$ of the endpoint table,
\begin{equation}
  \sum_P w_{\mathrm{full}}(P)F(L(P))
  =
  \sum_{L\ \mathrm{realized}}W(L)F(L).
  \label{eq:weighted-reference-grouping-v3}
\end{equation}
\end{proposition}

\begin{proof}
Partition the finite set of supports, with their physical decorations, by the
value of $L(P)$. On each fiber the factor $F(L(P))$ is constant. Summing first
inside each fiber gives $W(L)$, and then summing over the realized tables gives
the two identities.
\end{proof}

Combining Proposition~\ref{prop:fixed-support-all-deficit-v3},
\eqref{eq:support-cardinality-v3}, and
\eqref{eq:reference-grouping-v3} yields
\begin{equation}
  \sum_{\text{high skeletons}}w_{\mathrm{hi}}
  \le
  \left(\sum_{L\ \mathrm{realized}}W(L)\right)
  \left(1+(\alpha+1)\rho_{16}\right)^K.
  \label{eq:coarse-realized-table-reduction-v3}
\end{equation}

\begin{remark}[Sharper table-preserving form]
If every $\rho_{ij}\le1/2$, then
$\sum_{h\ge1}\rho_{ij}^h\le2\rho_{ij}$. Applying
\eqref{eq:weighted-reference-grouping-v3} with
$F(L)=\prod_{i,j}(1+2\rho_{ij})^{L_{ij}}$ gives
\[
  \sum_{\text{high skeletons}}w_{\mathrm{hi}}
  \le
  \sum_{L\ \mathrm{realized}}
    W(L)\prod_{i,j}(1+2\rho_{ij})^{L_{ij}}.
\]
The proof below uses the coarser form
\eqref{eq:coarse-realized-table-reduction-v3}, which requires only the single
condition $\rho_{16}\le1$ and cleanly separates the deficit sum from endpoint
transport.
\end{remark}

\begin{remark}[Reusable finite core]
The aggregate-fiber formula, the single ambient-denominator comparison, the
optional-choice identity, and the endpoint-table regrouping are not specific
to four endpoint sizes. They apply to any finite endpoint alphabet for which
the selected cells form a matching. The four consecutive sizes enter only in
the numerical deficit charge and the transport estimate below.
\end{remark}

\subsection{Endpoint transport}

Write the four endpoint sizes as
\[
  u_i=a-i,
  \qquad 0\le i\le3,
  \qquad a=\alpha-2.
\]
For a full-containment table $L=(\ell_{ij})$, set
\[
  r_i=\sum_j\ell_{ij},
  \qquad
  c_j=\sum_i\ell_{ij},
  \qquad
  x_{ij}=\min\{u_i,u_j\},
  \qquad
  J(L)=\sum_{i,j}x_{ij}\ell_{ij}.
\]
Its exact reference weight is
\begin{equation}
  W(L)=
  \frac{\prod_i(k_i)_{r_i}\prod_j(k_j)_{c_j}}
       {\prod_{i,j}\ell_{ij}!}
  \frac1{(n)_{J(L)}}
  \prod_{i,j}
  \left[
    \frac{(u_i)_{x_{ij}}(u_j)_{x_{ij}}}{x_{ij}!}g(x_{ij})
  \right]^{\ell_{ij}}.
  \label{eq:endpoint-reference-weight-v3}
\end{equation}
For a one-sided margin vector $r$, define
\begin{equation}
  D(r)=
  \frac{\prod_i(k_i)_{r_i}^2}{\prod_i r_i!}
  \frac{\prod_i[u_i!g(u_i)]^{r_i}}{(n)_{m_r}},
  \qquad
  m_r=\sum_i u_ir_i.
  \label{eq:one-sided-diagonal-weight-v3}
\end{equation}
This is the partial-diagonal weight from Section~7.

For the endpoint pair $(i,j)$, put
\[
  s_{ij}:=\min\{u_i,u_j\},
  \qquad
  t_{ij}:=\max\{u_i,u_j\},
  \qquad
  d_{ij}:=t_{ij}-s_{ij},
\]
and define
\[
  Q_{ij}:=
  (n+1)^{d_{ij}/2}
  \frac{\sqrt{(t_{ij})_{d_{ij}}}}{d_{ij}!}
  2^{-\{d_{ij}s_{ij}+\binom{d_{ij}}2\}/2}.
\]
Thus $Q_{ii}=1$. Since $t_{ij}\le a$, $s_{ij}\ge a-3$, and
$d_{ij}\le3$,
\begin{equation}
  Q_{ij}\le\frac{\eta_n^{d_{ij}}}{d_{ij}!},
  \qquad
  \eta_n=
  2^{3/2}\sqrt{\frac{(n+1)a}{2^a}}
  =O\!\left(\frac{(\log n)^{3/2}}{\sqrt n}\right).
  \label{eq:endpoint-eta-v3}
\end{equation}

Let
\[
  A_L=\frac{\prod_i r_i!}{\prod_{i,j}\ell_{ij}!},
  \qquad
  C_L=\frac{\prod_j c_j!}{\prod_{i,j}\ell_{ij}!},
  \qquad
  Q^L=\prod_{i,j}Q_{ij}^{\ell_{ij}}.
\]

\begin{lemma}[Square-free endpoint transport]
For every feasible endpoint table $L$,
\begin{equation}
  W(L)^2
  \le
  \bigl(D(r)A_LQ^L\bigr)\bigl(D(c)C_LQ^L\bigr).
  \label{eq:square-free-endpoint-transport-v3}
\end{equation}
\end{lemma}

\begin{proof}
The profile factors and the table factorials cancel exactly in the quotient of
$W(L)^2$ by $(D(r)A_L)(D(c)C_L)$. It remains to compare the ambient falling
factorials and the local full-containment atoms.

Fix one endpoint pair and abbreviate $s=s_{ij}$, $t=t_{ij}=s+d$, and
$x_{ij}=s$. Its local atom in \eqref{eq:endpoint-reference-weight-v3} is
\[
  \frac{(s)_s(t)_s}{s!}g(s)
  =
  \frac{t!}{d!}g(s).
\]
Consequently
\[
  \frac{
    \left[\frac{(s)_s(t)_s}{s!}g(s)\right]^2
  }{s!g(s)\,t!g(t)}
  =
  \frac{(t)_d}{(d!)^2}
  2^{-ds-\binom d2}
  =
  \frac{Q_{ij}^2}{(n+1)^d}.
\]
Here we used
$g(t)/g(s)=2^{ds+\binom d2}$; the identity is also valid for $d=0$.

Moreover, $m_r,m_c\ge J(L)$ and
\[
  m_r+m_c-2J(L)
  =
  \sum_{i,j}d_{ij}\ell_{ij}.
\]
Therefore
\[
  \frac{(n)_{m_r}(n)_{m_c}}{(n)_{J(L)}^2}
  \le
  (n+1)^{m_r+m_c-2J(L)}.
\]
This ambient factor cancels the denominators $(n+1)^{d_{ij}\ell_{ij}}$ from
the local ratios. Multiplying over all cells gives
\eqref{eq:square-free-endpoint-transport-v3}.
\end{proof}

\begin{proposition}[Endpoint-table sum]
\label{prop:endpoint-table-sum-v3}
Uniformly in the phase,
\begin{equation}
  \sum_LW(L)
  \le
  \exp\!\left\{O(\eta_nK)\right\}
  \sum_rD(r)
  \le
  \exp\!\left\{O(\sqrt{n\log n})\right\}.
  \label{eq:endpoint-table-sum-v3}
\end{equation}
The sum on the left may be enlarged from realized tables to all feasible
endpoint tables.
\end{proposition}

\begin{proof}
By the arithmetic--geometric mean inequality and
\eqref{eq:square-free-endpoint-transport-v3},
\[
  2W(L)
  \le
  D(r)A_LQ^L+D(c)C_LQ^L.
\]
Consider the first term. For a fixed row margin $r$, discard the column-margin
constraints; nonnegativity makes this an upper bound. The multinomial theorem
then gives
\[
  \sum_{L:\,\operatorname{row}(L)=r}A_LQ^L
  \le
  \prod_i\left(\sum_jQ_{ij}\right)^{r_i}
  \le
  (1+C\eta_n)^K.
\]
After multiplication by $D(r)$ and summation over $r$, the first half is at
most $(1+C\eta_n)^K\sum_rD(r)$. The second half is identical after exchanging
rows and columns. Section~7 proves $\sum_rD(r)=1+o(1)$ uniformly in the phase.
Finally, $K=\Theta(n/\log n)$ and \eqref{eq:endpoint-eta-v3} imply
$\eta_nK=O(\sqrt{n\log n})$.
\end{proof}

\subsection{Insertion of the phase estimates}

The phase satisfies
\[
  \alpha
  =
  \left(\frac{2}{\log 2}+o(1)\right)\log n.
\]
In particular, for all sufficiently large $n$,
\begin{equation}
  \frac52\log n\le\alpha.
  \label{eq:coarse-phase-corridor-v3}
\end{equation}
Every endpoint minimum satisfies $m_{ij}\ge\alpha-5$, and elementary floor
arithmetic gives
\[
  3\alpha-19
  \le
  4\left\lfloor\frac{3m_{ij}-1}{4}\right\rfloor.
\]
Using $\log 2>2/3$ and \eqref{eq:coarse-phase-corridor-v3}, we obtain
\begin{equation}
  \frac54\log n-\frac{19}{6}
  \le
  (\log 2)
  \left\lfloor\frac{3m_{ij}-1}{4}\right\rfloor.
  \label{eq:coarse-log-budget-v3}
\end{equation}
Hence
\[
  2^{\lfloor(3m_{ij}-1)/4\rfloor}
  \ge
  e^{-19/6}n^{5/4}.
\]
Since every $m_{ij}=O(\log n)$,
\begin{equation}
  \rho_{16}
  =O\!\left(\frac{\log n}{n^{1/4}}\right)
  \longrightarrow0.
  \label{eq:rho-sixteen-small-v3}
\end{equation}
Thus $\rho_{16}\le1$ for all sufficiently large $n$.

\begin{proposition}[Bare high-skeleton estimate]
\label{prop:bare-high-skeleton-v3}
There is a deterministic sequence $\varepsilon_n\to0$ such that
\begin{equation}
  \operatorname{BareSkeletonSum}_n
  \le
  \exp\!\left\{
    \varepsilon_n\frac{n}{(\log n)^4}
  \right\}.
  \label{eq:bare-skeleton-final-v3}
\end{equation}
\end{proposition}

\begin{proof}
The deficit factor in \eqref{eq:coarse-realized-table-reduction-v3} satisfies
\[
  \log\left(1+(\alpha+1)\rho_{16}\right)^K
  \le
  K(\alpha+1)\rho_{16}
  =O(n^{3/4}\log n).
\]
Proposition~\ref{prop:endpoint-table-sum-v3} contributes
$O(\sqrt{n\log n})$ to the logarithm. Both quantities are
$o(n/(\log n)^4)$. Combining these estimates with
\eqref{eq:coarse-realized-table-reduction-v3} proves the proposition.
\end{proof}

No residual local reward or binary cycle-space factor has been included in
$\operatorname{BareSkeletonSum}_n$. Section~9 treats exactly those remaining
factors.
